\begin{theorem}[The integral eight-manifold]\hypertarget{tag-02NW}{}\label{maps:src:F.2.5}
\StatementRecord{maps:src:F.2.5}{proved}{02NW}{\SeriesRef{maps:src:F.2.1}, \SeriesRef{night:P1.1}, \SeriesRef{night:P1.9}, \SeriesRef{night:P1.11}; \cite{ffAchingerOgus20,StacksProject,HirzebruchTopologicalMethods,ffPetersSterk}}


There is a closed smooth oriented eight-manifold with the prescribed
index-one Kuhn cusp, two free affine fillings of multiplicities $5,4$,
trivial fundamental group, and rational Betti numbers
\[
 (1,0,0,0,4,0,0,0,1). \tag{F.5}\hypertarget{tag-02F9}{}\label{maps:eq:F.5}
\]
It admits an almost complex structure.  Its signature is two and
its middle intersection lattice modulo torsion is
$3\langle1\rangle\oplus\langle-1\rangle$.
Its integral homology is computed in \cref{maps:src:F.2.6}, and a complex structure on it is constructed in \cref{maps:src:F.3.7}.

\end{theorem}

\begin{theorem}[The rank-six monodromy]\hypertarget{tag-02QI}{}\label{night:P1.1}
\StatementRecord{night:P1.1}{proved}{02QI}{the definitions and constructions stated in the proof}
The finite homology monodromies have orders $5,4$, respectively, and
characteristic polynomials
$(T-1)^2\Phi _5(T)$ and $(T-1)^2(T+1)^2(T^2+1)$.
For $N=M_0-I$ we have $N^2=0$, $\operatorname{rank}N=3$, and
nonzero Smith factors $1,1,1$.  Moreover
\[
 \Lambda^{\langle M_1,M_2\rangle}=0,\qquad
 \Lambda_{\langle M_1,M_2\rangle}=\mathbb Z,
\]
where the quotient homomorphism is $\ell(x)=x_0$.
On the cohomology lattice the corresponding groups are $\mathbb Ze_0$
and $\mathbb Z/2$.  Every nonidentity power of each specified affine
action $\tau_i$ is free.
\end{theorem}

\begin{theorem}[Integral homology and three-connectivity]\hypertarget{tag-02NX}{}\label{maps:src:F.2.6}
\StatementRecord{maps:src:F.2.6}{proved}{02NX}{\SeriesRef{maps:src:F.2.5}, \SeriesRef{night:P1.11}, \SeriesRef{night:P1.1}, \SeriesRef{maps:src:F.2.1}, \eqref{maps:eq:F.5}; \cite[Theorem 3.2, p.~195; Corollary 3.3, p.~196; Theorem 3.30, p.~241; Theorem 4.32, p.~366]{Hatcher2002}}


The manifold of Theorem \ref{maps:src:F.2.5} has free integral homology with
Betti numbers \eqref{maps:eq:F.5}.  In particular $H_2=H_3=0$, and it is
three-connected.

\end{theorem}

\begin{theorem}[The compact complex fourfolds $M_c$]\hypertarget{tag-02QJ}{}\label{night:P1.10}
\StatementRecord{night:P1.10}{proved}{02QJ}{\SeriesRef{night:P1.11}, \SeriesRef{night:P1.9}}
For every real parameter $c>-h_{\min}$ the period family $F_c$ of \cref{night:P1.8} gives a compact smooth complex fourfold $M_c$ and a proper holomorphic map
to $\mathbb P^1$ with the two multiple fibres of \cref{maps:src:F.2.5} and the
index-one Kuhn cusp.  With its complex orientation,
$\pi_1(M_c)=1$ and $\chi(M_c)=6$.
\end{theorem}

\begin{theorem}[The central cusp fibre and its collapse map]\hypertarget{tag-02QK}{}\label{night:P1.11}
\StatementRecord{night:P1.11}{proved}{02QK}{\SeriesRef{night:P1.9}}
The index-one central cusp fibre $C$ has torsion-free integral homology
with ranks $(1,3,9,9,10,3,1)$.
The collapse $\kappa:T^6\to C$ has primitive image in every degree;
it is onto except in degree four, where its cokernel is $\mathbb Z$.
Its image ranks are $(1,3,9,9,9,3,1)$.
In every degree its kernel is the saturation of
$(\bigwedge^kM_0-I)\bigwedge^k\Lambda$.
\end{theorem}

\begin{theorem}[Integral homology of $M_c$]\hypertarget{tag-02QL}{}\label{night:P1.12}
\StatementRecord{night:P1.12}{proved}{02QL}{\SeriesRef{night:P1.1}, \SeriesRef{night:P1.9}, \SeriesRef{night:P1.10}, \SeriesRef{night:P1.11}; \cite[Theorem 4.1.1 and Lemmas 4.1.2--4.1.3, pp.~480--482]{ffAchingerOgus20}; \cite[Theorem 3.2, p.~195; Theorem 3.30, p.~241; Theorem 4.32, p.~366]{Hatcher2002}}
The fourfold $M_c$ has torsion-free integral homology
\[
H_k(M_c;\mathbb Z)=
\begin{cases}
\mathbb Z,&k=0,8,\\
\mathbb Z^4,&k=4,\\
0,&k=1,2,3,5,6,7.
\end{cases}
\]
In particular it is three-connected.
\end{theorem}

\begin{theorem}[The genus-two pencil and its Hodge bundle]\hypertarget{tag-02QV}{}\label{night:P1.4}
\StatementRecord{night:P1.4}{proved}{02QV}{the definitions and constructions stated in the proof}
The genus-two pencil has one nodal fibre away from $0,\infty$, at
$\lambda_*=-256/3125$.  Good reduction at $0,\infty$ requires root
orders $8,5$.  On this root orbifold its Hodge bundle is
$\mathcal H=T^3\oplus T$, where
\[
T=\mathcal O\bigl((2/5)[\infty]-(3/8)[0]\bigr),\qquad
\deg T=1/40,
\]
and $T^3$ is the line generated meromorphically by $dx/y$.
\end{theorem}

\begin{theorem}[The rank-three bundle $P$]\hypertarget{tag-02QW}{}\label{night:P1.5}
\StatementRecord{night:P1.5}{proved}{02QW}{the definitions and constructions stated in the proof}
The wedge map from
$\ell_0\otimes(\mathcal H_{\mathrm{dR}}/\ell_0)$, followed by
projection modulo the nondegenerate polarization bivector, is
injective and has rank-three image. This image defines the bundle
$P$ used in the period construction.
\end{theorem}

\begin{theorem}[the geometric rank-five quotient]\hypertarget{tag-02O7}{}\label{maps:src:F.3.5}
\StatementRecord{maps:src:F.3.5}{proved}{02O7}{\SeriesRef{maps:src:F.2.5}}


Put $\lambda_*=-256/3125$ and
$U=\mathbb P^1\setminus\{0,\infty,\lambda_*\}$.  The smooth projective
models of
\[
 C_\lambda:\quad y^2=x^5+\lambda^3x+\lambda^4
 \tag{F.7}\hypertarget{tag-02FB}{}\label{maps:eq:F.7}
\]
give a polarized rank-four rational cohomology system $\mathcal H$ on
$U$.  Its primitive exterior square is rationally isomorphic to
the rank-five quotient of the system with matrices $A,B$ in
Theorem \ref{maps:src:F.2.5} by its invariant first coordinate line.
On the orbifold $\mathcal B'$ obtained by taking roots of orders
five at infinity and eight at zero, the extended Hodge bundle is
\[
 \mathcal E=\mathcal T^3\oplus\mathcal T,\qquad
 \mathcal T=\mathcal O_{\mathcal B'}(2D_\infty-3D_0),\qquad
 \deg\mathcal T=\frac1{40}.
 \tag{F.8}\hypertarget{tag-02FC}{}\label{maps:eq:F.8}
\]
Here $D_\infty,D_0$ are the reduced root divisors; the extension at
$\lambda_*$ is the nilpotent canonical extension.

\end{theorem}

\begin{theorem}[The integral exterior-square quotient]\hypertarget{tag-02QT}{}\label{night:P1.2}
\StatementRecord{night:P1.2}{proved}{02QT}{the definitions and constructions stated in the proof}
The lower blocks of $A,B$ are precisely the quotient representations of
$\bigwedge^2(C_f^{-1}),\bigwedge^2C_g$ in the basis $Q$.
The integral quotient is the dual of the primitive exterior-square
lattice.  The natural map from the primitive lattice into that quotient
has cokernel $\mathbb Z/2$; these two integral representations are not
isomorphic.  Their rationalizations are isomorphic by that same map.
\end{theorem}

\begin{theorem}[Integral marking of the pencil]\hypertarget{tag-02QY}{}\label{night:P1.7}
\StatementRecord{night:P1.7}{proved}{02QY}{\SeriesRef{maps:src:F.3.5}, \SeriesRef{night:P1.2}}
The integral $H^1$ monodromy of the pencil has a companion marking with
finite pair $(C_f^{-1},C_g)$, and cusp product $C_fC_g^{-1}$.
Its exterior-square quotient is consequently the integral quotient
representation appearing in \cref{night:P1.2}.
\end{theorem}

\begin{theorem}[a global nonpolarized period family]\hypertarget{tag-02O8}{}\label{maps:src:F.3.6}
\StatementRecord{maps:src:F.3.6}{proved}{02O8}{\eqref{maps:eq:F.11}, \eqref{maps:eq:F.14}, \eqref{maps:eq:F.16}, \eqref{maps:eq:F.19}, \eqref{maps:eq:F.12}, \SeriesRef{maps:src:F.2.1}, \eqref{maps:eq:F.22}, \eqref{maps:eq:F.20}, \eqref{maps:eq:F.21}}


The integral rank-six system with cohomology monodromies $A,B$
admits a holomorphic subbundle $\mathcal F$ of rank three over
$U$ satisfying $\mathcal F\cap\overline{\mathcal F}=0$.
It extends holomorphically on the finite root charts.  At the
cusp an integral marking puts its periods in the form
\[
 Z(t)=\frac{\log t}{2\pi i}I_3+C(t),
 \tag{F.13}\hypertarget{tag-02EJ}{}\label{maps:eq:F.13}
\]
where $C$ is holomorphic at zero and bounded there.

\end{theorem}

\begin{theorem}[The exact parameter domain of the lifted family]\hypertarget{tag-02QZ}{}\label{night:P1.8}
\StatementRecord{night:P1.8}{proved}{02QZ}{\SeriesRef{night:P1.5}}
With the rank-three bundle $P$ of \cref{night:P1.5}, the rank-six
system admits a holomorphic lifted family
$F_c=(\sigma+ic\alpha e)(P)$ on the punctured base.
For a specified splitting $\sigma$ there is a real analytic function
$h$ with finite attained minimum $h_{\min}$, and the exact domain of
real parameters giving $F_c\cap\overline{F_c}=0$ on the completed
base, including both finite root points and excluding only the cusp
before its toroidal filling, is $c>-h_{\min}$.
At the cusp its periods have the form
$Z(q)=(\log q/(2\pi i))I_3+C(q)$ with $C$ holomorphic and bounded.
The construction extends to the two finite root charts and descends to
orders $5,4$ on the rank-six system.
\end{theorem}

\begin{theorem}[The Kuhn filling at the cusp]\hypertarget{tag-02R0}{}\label{night:P1.9}
\StatementRecord{night:P1.9}{proved}{02R0}{the definitions and constructions stated in the proof}
The periodic Kuhn fan, in the marking $S$, gives a smooth proper
holomorphic filling of the periods $Z(q)=sI_3+C(q)$ for a sufficiently
small disk.  Its central fibre is reduced with normal crossings.
Modulo the index-one period lattice it has respectively $1,7,12,6$
strata of dimensions $3,2,1,0$, each isomorphic to the corresponding
complex torus $(\mathbb C^*)^d$, and Euler number $6$.
An index-$d$ period sublattice gives $d$ times these stratum counts and
Euler number $6d$.
\end{theorem}

\begin{theorem}[compact holomorphic filling]\hypertarget{tag-02O9}{}\label{maps:src:F.3.7}
\StatementRecord{maps:src:F.3.7}{proved}{02O9}{\SeriesRef{maps:src:F.2.6}, \SeriesRef{maps:src:F.2.5}, \eqref{maps:eq:F.24}, \eqref{maps:eq:F.13}, \eqref{maps:eq:E.4}, \eqref{maps:eq:E.5}, \SeriesRef{maps:src:E.1.2}, \SeriesRef{maps:src:E.1.3}}


The smooth eight-manifold of Theorem \ref{maps:src:F.2.6} carries an integrable
complex structure.  There is a proper holomorphic map to
$\mathbb P^1$ whose regular fibres are complex three-tori,
whose two finite reduced fibres are the free affine quotients
of orders five and four specified in Theorem \ref{maps:src:F.2.5}, and whose
remaining fibre is the index-one Kuhn degeneration.

\end{theorem}

\begin{theorem}[Characteristic numbers and the middle lattice]\hypertarget{tag-02QM}{}\label{night:P1.13}
\StatementRecord{night:P1.13}{proved}{02QM}{\SeriesRef{night:P1.12}, \SeriesRef{sec:V-known-inputs}; \cite[equations (1.2)--(1.4), pp.~3--4]{ffEsoleKang}; \cite[Theorem 2.2.1, p.~61]{ffPetersSterk}}
With the complex orientation the fourfold has
\[
\sigma(M_c)=2,\qquad H^4(M_c;\mathbb Z)\cong I_{3,1},\qquad
w_2(M_c)=0,
\]
and there is a diagonal integral basis in which
$p_1/2=(1,1,1,1)$.
Its Chern numbers are
$c_4=6$, $c_2^2=2$, and
$c_1^4=c_1^2c_2=c_1c_3=0$; its Pontryagin numbers are
$p_1^2=8$, $p_2=14$.
\end{theorem}

\begin{theorem}[a complex structure on the stated connected sum]\hypertarget{tag-02OC}{}\label{maps:src:F.4.1}
\StatementRecord{maps:src:F.4.1}{proved}{02OC}{\SeriesRef{maps:src:F.3.7}, \SeriesRef{maps:src:F.2.6}, \SeriesRef{sec:V-known-inputs}, \SeriesRef{night:P1.13}; \cite[Corollary 1.4, p.~4]{ffMandal}; \cite[Definition 3.6 and Theorem 3.7(1), pp.~12--14]{IVVWallForm}; \cite[Theorem A, p.~178]{IVVKramerStolz}}


The compact complex fourfold of Theorem \ref{maps:src:F.3.7} is three-connected,
has Euler number six, and is orientation-preservingly diffeomorphic to
\[
 3\mathbb H\mathbb P^2\ \#\ (-\mathbb H\mathbb P^2).
 \tag{F.25}\hypertarget{tag-02EW}{}\label{maps:eq:F.25}
\]
Its integral middle lattice is $I_{3,1}$ and, in a diagonal integral
basis, $p_1/2=(1,1,1,1)$.  

\end{theorem}